\begin{RequAnswer}{3.1}
(b), (d), (f) make sense.
For the incorrect ones, first note that $A$ and $B$ are both sets of numbers.
(a) does not make sense because an element cannot be a subset of a set.
For (c), $A$ contains the number $3$, but not the set containing $3$,
which is what $\{3\}$ is.
For (e), $A$ does not contain as an element another set (i.e. $B$), so that notation does not
make sense.
\end{RequAnswer}
\begin{RequAnswer}{3.2}
5. Remember, repeated elements do not count!
\end{RequAnswer}
\begin{RequAnswer}{3.3}
(a) i. Yes, ii. No, iii. No.
(b) 6, 3, 4.
(c) They are all finite sets.
\end{RequAnswer}
\begin{RequAnswer}{3.4}
(a) Yes
(b) No
(c) Yes
(d) Yes
(e) No
\end{RequAnswer}
\begin{RequAnswer}{3.5}
$\{n^3 | n\in \BBZ\}$
\end{RequAnswer}
\begin{RequAnswer}{3.6}
$\{n/100 | n\in \BBZ\}$
\end{RequAnswer}
\begin{RequAnswer}{3.7}
Yes (because the power set of $A$ is the set of all subsets of $A$, and $A\subseteq A$, so $A\in P(A)$)
and No (because the elements of $P(A)$ are subsets of elements of $A$, so a subset of $P(A)$ would
be a set of subsets of $A$, but $A$ is a set of elements (of whatever type $A$ consists of).
This is a subtle point, so do not worry too much about it unless you plan to major in mathematics.
\end{RequAnswer}
\begin{RequAnswer}{3.8}
$2^5=32$.
\end{RequAnswer}
\begin{RequAnswer}{3.9}
(a) $|A|=5$ and $|P(a)|=2^5=32$.
(b) No. The elements of $A$ are $a$, $b$, etc.
If that statement were true, $\{a\}\in A$, but it is not. Remember, $a\in A$ (which is true) and
$\{a\}\in A$ (which is not true) are not saying the same thing.
(c) Yes.
(d) No. As in part (b), the elements of $A$ are $a$, $b$, etc. but $\{b,c,e\}$ is a set of those elements.
Remember: $\in$ means ``element of,'' and $\subseteq$ means ``subset of,'' which are not the same thing!
(e) Yes because each of the three sets in the set on the left are subsets of $A$.
(f) No. To be a subset of the power set, it needs to be a set of subsets. This is a set of elements.
(g) Yes because the set $\{b,c,e\}\subseteq A$.
\end{RequAnswer}
\begin{RequAnswer}{3.10}
(a) $\{a, b, c, d, f, g, h, j, k, l, m, n, p, q, r, s, t, v, w, x, y, z\}$
(b) $\{b, d, g, k, p, v\}$
(c) $\{a\}$ (NOT $a$!)
(d) $\{a\}$
(e) $\{e, i, o, u \}$
\end{RequAnswer}
\begin{RequAnswer}{3.11}
(a) True;
(b) False;
(c) False;
(d) True (This is DeMorgan's law in words)
\end{RequAnswer}
\begin{RequAnswer}{3.12}
(a) answers will vary, but should be things like $(1,z)$, $(3,x)$, and $(2,v)$.
(b) answers will vary, but should be things like $(1,2)$, $(3,3)$, and $(4,1)$.
(c) 24.
(d) $4^3=64$.
(e) $2^{4^2*6}=2^{96}$.
\end{RequAnswer}
\begin{RequAnswer}{3.13}
The area being pointed to is ``with {\bf AND} without you,''
so it is not at all correct.

\noindent
\begin{minipage}[t]{.62\textwidth}
%\vspace{0pt}
That is definitely not where Bono can't live.
Notice that ``with you'' and ``without you'' are complements,
so the Venn diagram to the right demonstrates a proper relationship
between them. So where can't Bono live? He can't live anywhere in the rectangle (the universe)
because he can't live in the circle (``with you'') or outside the circle (``without you'').
Put another way, the entire diagram is the place where Bono can't live.
\end{minipage}
\hfill
\begin{minipage}[t]{.27\textwidth}
\vspace{0pt}
\psset{unit=1.25pc}
%\begin{pspicture}(-3, -.25)(3,4.1)
\begin{pspicture}(3, 2)(-3.5,-3)
\pscustom[fillstyle=solid,fillcolor=lightgray]{\psline(-3.5,2.5)(-3.5,-2.5)(3.5,-2.5)(3.5,2.5)(-3.5,2.5)}
\uput[l](-.25,1.75){without}
\uput[l](-1,.75){you}
\pscircle[fillstyle=solid,fillcolor=white](1,0){2}
\uput[u](1,0){with}
\uput[u](1,-1){you}
%\uput[r](4,2){$U$}
\uput[r](-3.75,-3){Where Bono can't live}
\end{pspicture}
%\meinecaption{1}{$\overline{A}$}
%\label{fig:a_complement}
\end{minipage}
\end{RequAnswer}
\begin{RequAnswer}{3.14}
(a) $A\subseteq B$.
(b) $A=B$.
\end{RequAnswer}
\begin{RequAnswer}{3.15}
First, $A\cap B$ is the set of elements that are in both $A$ and $B$. Therefore,
$\overline{A\cap B}$  is the set of elements that are not in both $A$ and $B$.
(This includes elements that are in $A$ but not $B$, in $B$ but not $A$, or in neither $A$ nor $B$.)
$\overline{A}\cup\overline{B}$ is the union of the elements that are not in $A$ and the elements that are not in $B$.
That is, it is any element that is either not in $A$ or not in $B$. So it is all elements that are not in both $A$ and $B$,
which is exactly what $\overline{A\cap B}$ is.
Therefore, $\overline{A\cap B}= \overline{A}\cup\overline{B}$.
\end{RequAnswer}
\begin{RequAnswer}{3.16}
First, notice that since $U$ is the universal set, $A\subseteq U$ and $\overline{A}\subseteq U$.
Let $x\in A\cup \overline{A}$. Then either $x\in A$ or $x\in\overline{A}$.
In either case, $x\in U$ (since $A\subseteq U$ and $\overline{A}\subseteq U$). Therefore $A\cup \overline{A}\subseteq U$.

Now, assume $x\in U$.  If $x\in A$, then $x\in A\cup \overline{A}$.
If $x\not\in A$, then by definition, $x\in\overline{A}$, and therefore $x\in A\cup \overline{A}$.
In either case, $x\in A\cup \overline{A}$. Therefore $U\subseteq A\cup \overline{A}$.

Since we proved containment both ways, $A\cup \overline{A}=U$.
\end{RequAnswer}
\begin{RequAnswer}{3.17}
Given an integer $a$, compute $a\bmod 2$.
If it is $0$, $a$ is even, and if it is $1$, $a$ is odd.
\end{RequAnswer}
\begin{RequAnswer}{3.18}
It certainly can be used. If the current time is $t$, and you are on 24-hour time,
then in 8 hours the time will be $(t+8 \bmod 24$. If you are using 12-hour time, it is a little more complicated.
You can compute $(t+8)\bmod 12$ to get the time in 8 hours, but if the result is $0$,
you need to treat it like 12.
Also, this does not take into account whether there was a switch from am to pm.
That is more complicated and we would need to use additional logic to get that part correct.
\end{RequAnswer}
\begin{RequAnswer}{3.19}
Compute $c=a\bmod n$ and $d=b\bmod n$.
According to Theorem~\ref{thm:congruence},
$c=d$ if and only if $a\equiv b\pmod n$.
so just determine whether or not $c=d$.
\end{RequAnswer}
\begin{RequAnswer}{3.20}
$\gcd(524,118)=15$ as demonstrated here:

\begin{tabular}{|l|l|l|l|}
\hline
step&a&b&r\\
\hline
\hline
1&67890&12345&6165\\
\hline
2&12345&6165&15\\
\hline
3&6165&15&0\\
\hline
4&15&0&0\\
\hline
\end{tabular}
\end{RequAnswer}
\begin{RequAnswer}{3.21}
They are not because they have a common factor of 15.
\end{RequAnswer}
\begin{RequAnswer}{3.22}
No. In fact, they are usually not going to be the same (unless you start with an integer).
As an example, the floor of 3.2 is 3, and the ceiling of 3 is 3, so the ceiling of the floor of 3.2 is 3.
The ceiling of 3.2 is 4, and the floor of 4 is 4, so the floor of the ceiling of 3.2 is 4.
\end{RequAnswer}
\begin{RequAnswer}{3.23}
(a) $\BBZ$. (b) $\BBZ$. (c) $\{2z | z\in\BBZ\}$. That is, the set of all even numbers.
\end{RequAnswer}
\begin{RequAnswer}{3.24}
Here are possible answers. There are many other possible correct answers.
(a) $f(x)=2x$.
(b) $f(x)=6$.
(c) $f(x)=2x$.
(d) $f(x)=4$.
(e) $f(x)=2x$.
\end{RequAnswer}
\begin{RequAnswer}{3.25}
Here are possible answers. There are many other possible correct answers.
(a) $f(x)=2x$.
(b) $f(x)=(2x)\bmod 6$ (so $f(4)=2$ instead of 8).
(c) Impossible. The domain only has 4 numbers, and the codomain has 6, so it is impossible to map to all of them.
(d) $f(x)=2x$ (it doesn't map to 10 or 12).
(e) Impossible. Since an onto function is impossible, a bijective one is as well.
\end{RequAnswer}
\begin{RequAnswer}{3.26}
$(f\circ g)(x)=2^{x+2}$ and $(g\circ f)(x)=2^x+2$.
Did you get them backwards? If so, look at the definition of composition of functions again.
\end{RequAnswer}
\begin{RequAnswer}{3.27}
(a) False. If $a=b$, then clearly $f(a)=f(b)$, so that's meaningless. You need to show that if $f(a)=f(b)$, then $a=b$.
(b) True.
(c) False. That's the definition of a function. to show onto, you need to show that every element of the codomain
gets mapped to by some element of the domain.
(d) True. Since the range is the set of values actually mapped to, if it equals the codomain, then every value is mapped
to and the function is onto.
(e) False. For two reasons. First, elements of the range are always mapped to--that's the definition of range.
Even if you change range to codomain, it is still false. You only need to show that there is at least one element of the
codomain that is not mapped to.
(f)True.
\end{RequAnswer}
\begin{RequAnswer}{3.28}
First, notice that if $f(a)=f(b)$, then $a^3-8=b^3-8$,
which implies $a^3=b^3$. Taking the third root of each side,
we get $a=b$. Thus, $f$ is one-to-one.
Now, notice that if $x\in\reals$, then $\sqrt[3]{x+8}\in \reals$,
and $f\left(\sqrt[3]{x+8}\right)=\left(\sqrt[3]{x+8}\right)^3-8=x+8-8=x$,
so $f$ is onto.
To find the inverse (which I already did in order to prove it was
onto, but I'll show the work anyway),
we solve $y=x^3-8$ for $x$. We get $x^3=y+8$, so
$x=\sqrt[3]{y+8}$. Replacing variables, we have
that $f^{-1}=\sqrt[3]{x+8}$.
\end{RequAnswer}
\begin{RequAnswer}{3.29}
(a) $C=\{1,3,5,7,9\}$.
(b) $C=\{1,2,3,5,7,9\}$ (or any set that contains all of 1,3,5,7,9, and at least one of 2, 4, 6, 8, or 10).
(c) $C=\{1,3,7,9\}$ (or any proper subset of $\{1,3,5,7,9\}$).
(d) $C=\{1,3,5\}$ and $D=\{7,9\}$ (or any two disjoint sets such that $C\cup D=\{1,3,5,7,9\}$.
(e) $C=\{1,2,3,5\}$ and $D=\{7,9\}$ (or many other possibilities).
\end{RequAnswer}
\begin{RequAnswer}{3.30}
Yes. It is a subset of $\BBZ\times\BBZ$.
\end{RequAnswer}
\begin{RequAnswer}{3.31}
There are many possible answers. Here is one.
Define $T$ to be the relation on human beings such that $x T y$ if and only if $x$ is at least as tall as $y$.
It is not hard to see that $T$ is reflexive, anti-symmetric, and transitive, so it is a partial order.
\end{RequAnswer}
\begin{RequAnswer}{3.32}
answers we vary, but here is one possibility.
Let $C$ be the set of all cars.
Define the relation $R$ on $C$ by $x R y$ if and only if $x$ was manufactured in the same year as $y$.
It is not hard to see that $R$ is symmetric, reflexive, and transitive, so it is an equivalence relation.
Let $C_y$ be the set of all cars manufactured in year $y$.
Then it is not hard to see that $C=C_{1886}\cup C_{1887}\cup\cdots\cup C_{2023}\cup C_{2024}\cup C_{2025}$
is a partition of $C$ (assuming you aren't reading this in 2026 or later, and assuming we regard 1886 as the first year
a car was made, which turns out to be a question without a clear answer).
For a representative for class $C_y$, simply pick any car that was manufactured in year $y$.
\end{RequAnswer}
\begin{RequAnswer}{3.33}
(a) It is straightforward to prove that $B$ is symmetric, reflexive, and transitive, so it is an equivalence relation.
(b) It is not. It is no anti-symmetric.
(c) Let $B_i$ be the positive integers that have $i$ ones in their binary representation.
It is easy to see that $B_i\cap B_j=\emptyset$ if $i\not = j$ and that  $\BBZ^+ = B_1\cup B_2\cup B_3\cup\cdots$,
so this definition gives us a partition of $\BBZ^+$.
(d) We just let $a_i$ be the number whose binary representation is $i$ 1s.
For instance, $a_1=1$, $a_2=3$, $a_3=7$, etc. Notice that we can do better by defining it explicitly: $a_i=2^i-1$.
(e) The smallest elements of $B_2$ would be $3=11_2$, $5=101_2$, $6=110_2$, and $9=1001_2$.
There are infinitely many other possible answers.
\end{RequAnswer}
